\begin{lemma}[A nonzero multi-index has order at least one]
Fix $n \in \mathbb{N}_0$ and let $\alpha \in \mathbb{N}_0^n$ be a multi-index
with order $|\alpha| = \sum_{j=1}^n \alpha_j$ as in \noderef{MultiIndexSum}.
If $\alpha \neq 0$, that is, if $\alpha$ is not the multi-index all of whose
entries vanish, then
\[
  1 \leq |\alpha| ,
\]
an inequality between real numbers: the order $|\alpha|$ is a natural number,
and it is its image under the canonical cast $\mathbb{N}_0 \to \mathbb{R}$ that
is compared with the real number $1$.
\end{lemma}
\begin{proof}
Step 1 (the order is a nonzero natural number).  We claim $|\alpha| \neq 0$ as a
natural number.  Suppose instead that $|\alpha| = 0$, that is
$\sum_{j=1}^n \alpha_j = 0$.  Every summand $\alpha_j$ is a natural number and
so is nonnegative, and a finite sum of nonnegative integers vanishes only if
every one of its terms vanishes; hence $\alpha_j = 0$ for every index
$j \in \{1,\ldots,n\}$.  Two multi-indices that agree at every index are equal,
so $\alpha = 0$, contradicting the hypothesis $\alpha \neq 0$.  Therefore
$|\alpha| \neq 0$.

Step 2 (a nonzero natural number is at least one).  In $\mathbb{N}_0$ the
conditions $m \neq 0$ and $1 \leq m$ are equivalent, since the natural numbers
below $1$ are exactly $\{0\}$.  Applying this to $m = |\alpha|$ and using
Step 1 gives
\[
  1 \leq |\alpha| \qquad \text{in } \mathbb{N}_0 .
\]

Step 3 (transfer to the reals).  The canonical cast $\mathbb{N}_0 \to \mathbb{R}$
is an order embedding: for natural numbers $m$ and $m'$ one has $m \leq m'$ if
and only if the corresponding reals satisfy the same inequality, and the cast
sends $1$ to the real number $1$.  Applying the forward direction to the
inequality of Step 2 therefore yields $1 \leq |\alpha|$ as an inequality in
$\mathbb{R}$, which is the assertion.  This cast is the only place the statement
mixes the two number systems: the hypothesis and Steps 1--2 are statements about
natural numbers, and only the conclusion is read in $\mathbb{R}$.
\end{proof}
